\documentclass[11pt,a4paper]{article}
\usepackage[utf8]{inputenc}
\usepackage[T1]{fontenc}
\usepackage{lmodern}
\usepackage[margin=1in]{geometry}
\usepackage{enumitem}
\usepackage{xcolor}
\usepackage{hyperref}
\usepackage{booktabs}
\usepackage{longtable}
\usepackage{array}
\usepackage[most]{tcolorbox}

% Define custom colors
\definecolor{promptblue}{RGB}{41, 98, 150}
\definecolor{promptbg}{RGB}{240, 247, 255}
\definecolor{promptborder}{RGB}{100, 160, 220}
\definecolor{accentgreen}{RGB}{0, 128, 100}
\definecolor{warningorange}{RGB}{200, 120, 0}

\newtcolorbox{resultbox}[1][]{
    colback=promptbg,
    colframe=promptborder,
    coltitle=white,
    fonttitle=\bfseries\small,
    title={#1},
    colbacktitle=promptblue,
    boxrule=1pt,
    arc=2mm,
    left=8pt, right=8pt, top=6pt, bottom=6pt,
    enhanced, breakable
}

\newtcolorbox{gapbox}[1][]{
    colback=orange!5,
    colframe=warningorange,
    coltitle=white,
    fonttitle=\bfseries\small,
    title={#1},
    colbacktitle=warningorange,
    boxrule=1pt,
    arc=2mm,
    left=8pt, right=8pt, top=6pt, bottom=6pt,
    enhanced, breakable
}

\hypersetup{
    colorlinks=true,
    linkcolor=promptblue,
    urlcolor=promptblue,
    citecolor=promptblue
}

\title{\textbf{Cross-Study Synthesis: machine learning }}
\author{Generated by Research Accelerant Agent}
\date{\today}

\begin{document}
\maketitle


\textbf{\large Studies Included in Synthesis}\par\vspace{0.5em}

\begin{longtable}{p{5cm}cc}
\toprule
\textbf{Study} & \textbf{Year} & \textbf{Citations} \\
\midrule
\endhead
A Survey on IoT Positioning Leveraging LPWAN, GNSS... & 2023 & 104 \\
Maritime Autonomous Surface Ships: A Review of Cyb... & 2024 & 48 \\
Robust Automatic Modulation Classification in Low ... & 2023 & 47 \\
Unmanned Autonomous Intelligent System in 6G Non-T... & 2024 & 35 \\
Machine Learning and Deep Learning powered satelli... & 2023 & 34 \\
\bottomrule
\end{longtable}

\vspace{1em}


\textbf{\large 1. Methodological Patterns}\par\vspace{0.5em}

Across the 5 reviewed studies on machine learning , the following methodological patterns emerged:

- **Computational / Simulation**: 2 studies
- **Cross-Sectional Survey**: 1 study
- **Not specified**: 1 study
- **Qualitative Case Study**: 1 study

**Design Distribution:**
- Total experimental studies: 0
- Total meta-analyses/systematic reviews: 0
- Total longitudinal studies: 0
- Total qualitative studies: 1
- Total computational studies: 2

\vspace{1em}


\textbf{\large 2. Overarching Findings and Trends}\par\vspace{0.5em}

**Consensus Themes:**
- Multiple studies focus on: autonomous, challenges, future, learning.

**Emergent Trends:**
- 5 recent studies (2023-2023) indicate evolving methodological approaches.

**Contradictions/Debates:**
- No major contradictions noted.

\vspace{1em}


\textbf{\large 3. Recurring Gaps and Limitations}\par\vspace{0.5em}

**Limitations noted by authors:**
In "Robust Automatic Modulation Classification in Low Signal to Noise Ratio" (2023), authors note: limited.

**Systematic gaps in the literature:**
1. Geographic scope: Most studies on machine learning  may be concentrated in specific regions, limiting generalizability to diverse populations.
2. Cross-method integration: While diverse methods exist (Cross-Sectional Survey, Not specified, Computational / Simulation, Qualitative Case Study), few studies combine approaches for robust validation.
3. Sample and temporal constraints: Many studies use cross-sectional designs or limited sample sizes, limiting causal inference about machine learning .
4. Translation to practice: Limited evidence on how findings regarding machine learning  translate into actionable interventions or policy recommendations.

\vspace{0.5em}


\begin{gapbox}[Gap 1]
Geographic scope: Most studies on machine learning  may be concentrated in specific regions, limiting generalizability to diverse populations.
\end{gapbox}

\vspace{0.3em}


\begin{gapbox}[Gap 2]
Cross-method integration: While diverse methods exist (Cross-Sectional Survey, Not specified, Computational / Simulation, Qualitative Case Study), few studies combine approaches for robust validation.
\end{gapbox}

\vspace{0.3em}


\begin{gapbox}[Gap 3]
Sample and temporal constraints: Many studies use cross-sectional designs or limited sample sizes, limiting causal inference about machine learning .
\end{gapbox}

\vspace{0.3em}


\begin{gapbox}[Gap 4]
Translation to practice: Limited evidence on how findings regarding machine learning  translate into actionable interventions or policy recommendations.
\end{gapbox}

\vspace{0.3em}


\vspace{1em}
\textbf{\large 4. Impact Assessment}\par\vspace{0.5em}

**High-Impact Studies (>$50$ citations):** 1 out of 5 reviewed studies (20%).

**Recent Contributions (2024+):** 2 studies.

**Most Cited Study:** A Survey on IoT Positioning Leveraging LPWAN, GNSS, and LEO-PNT (104 citations).

**Open Access Availability:** 5 of 5 studies (100%) provide open access to their full text.

\vspace{1em}


\textbf{\large 5. Future Research Directions}\par\vspace{0.5em}

Based on the identified gaps, the following research directions are recommended:

1. **Geographic scope**: Further investigation needed to address geographic scope: most studies on machine learning  may be concentrated in specific regions, limiting generalizability to diverse populations..
2. **Cross-method integration**: Further investigation needed to address cross-method integration: while diverse methods exist (cross-sectional survey, not specified, computational / simulation, qualitative case study), few studies combine approaches for robust validation..
3. **Sample and temporal constraints**: Further investigation needed to address sample and temporal constraints: many studies use cross-sectional designs or limited sample sizes, limiting causal inference about machine learning ..
4. **Translation to practice**: Further investigation needed to address translation to practice: limited evidence on how findings regarding machine learning  translate into actionable interventions or policy recommendations..

\vspace{1em}


\vfill
\begin{center}
\textit{Synthesis generated by Research Accelerant Agent on \today.}
\end{center}

\end{document}